\chapter{Conclusion}

This report has consolidated every single-aperture optimization model this project has
built: an exact wave-optics blur-minimization pipeline, a brightness-constrained
variant of it, a full human-vision acuity layer built on standard photometric and
contrast-sensitivity theory, and a structurally independent ideal-observer Monte Carlo
model built specifically to cross-check the others. It has also consolidated this
project's own audit of itself --- a real, confirmed optical bug, independently
re-derived two ways and fixed, whose correction moved the headline optimum-aperture
number by as much as $150\%$ at some myopia levels; and a separate import-shadowing
bug that meant two of the project's own acuity-validation efforts had, for a period,
never actually tested the pipeline they were built to validate.

The physics is internally consistent and, where checkable against closed-form theory
(the zero-defocus Airy disk, the Fresnel on-axis oscillation formula), matches it
closely. Two independent, physically-motivated corrections --- the wavelength fix and a
brightness-admission floor --- each move the naive blur-only optimum substantially
toward the externally observed clinical bracket, without either fully explaining it
alone. Where this project's models are tested against real clinical acuity data,
however, they do not yet pass: this is reported plainly, with the specific numbers,
rather than reframed as a partial success. The most defensible, load-bearing
conclusion this report can draw is therefore not a single aperture number, but a map
of exactly which models predict what, under which explicit assumptions, and precisely
where each one is --- and is not --- currently validated. Table~\ref{tab:synthesis}
(Chapter~\ref{ch:synthesis}) is that map.

Closing this project's remaining gap between theory and clinical fact, per
Chapter~\ref{ch:limitations}, does not appear to require a fifth acuity model or a
more cleverly parametrized correction term. It requires new, real, per-subject
clinical data, and eventually a physical prototype tested against it --- the one step
this entire body of theoretical work has been building toward and has not yet taken.
